\documentclass[11pt,a4paper]{article}
\usepackage[margin=2.5cm]{geometry}
\usepackage{amsmath,amssymb,amsthm,mathtools,booktabs,array,xcolor}
\usepackage[colorlinks=true,linkcolor=blue!60!black,citecolor=green!40!black,urlcolor=blue!60!black]{hyperref}
\newcommand{\La}{\textsf{\textbf{(a)}}}
\newcommand{\Lb}{\textsf{\textbf{(b)}}}
\newcommand{\Lc}{\textsf{\textbf{(c)}}}
\newcommand{\Ld}{\textsf{\textbf{(d)}}}
\newcommand{\script}[1]{\texttt{scripts/#1}}
\newtheorem{proposition}{Proposition}[section]
\theoremstyle{remark}\newtheorem{remark}[proposition]{Remark}

\title{Towards quantisation of Boussinesq solitons:\\ Phase 2 master report (Routes A--E, Phases 2a--2h)}
\author{Project \emph{quantum\_boussinesq\_solitons} --- PI: Gustav Delius}
\date{Updated 4 November 2026}
\begin{document}\maketitle

\begin{abstract}
Following the PI's initial decision, we first studied the \emph{bad} Boussinesq equation
\[
v_{tt}=v_{xx}+v_{xxxx}+3(v^2)_{xx},
\]
whose soliton sector is regular and purely elastic, has regular breathers, and is the continuum limit of the quantum Toda chain (Phase~1), and subsequently extended the analysis to the \emph{good} (well-posed) Boussinesq equation $u_{tt}=u_{xx}-u_{xxxx}-3(u^2)_{xx}$ and its Corrigan--Zambon integrable defects (Phases~2c--2h).

\paragraph{Semiclassical ($O(\hbar^{-1})$) results (Route A).}
Two-soliton scattering is an exactly symplectic map for the canonical structure of the field theory, \emph{only} once ordering-dependent ``topological'' cross terms $d\Pi_i\wedge dM_j$ are included. These terms come from the kink in the potential $\phi$. With them we obtain a closed-form semiclassical two-soliton phase, $S\simeq e^{i\delta/\hbar}$:
\[
\delta(1,2)=\tfrac29(\ell_1-\ell_2)\log A_{12}+\tfrac23M_1M_2(v_1-v_2),\qquad \ell=\epsilon\cosh3\beta\ \text{(bad)},\quad \ell=-\epsilon\cos3\alpha\ \text{(good)}.
\]
\begin{itemize}
\item $\delta$ is antisymmetric (unitarity) and obeys a crossing-type relation under $\beta\to\beta+i\pi$.
\item Continued to complex conjugate rapidities, $i\cdot\tfrac29(\ell_1-\ell_2)$ reproduces \emph{exactly} the adiabatic invariant of the Boussinesq breather computed independently by Bogdan and Charkina \cite{BC}, turning breather quantisation into a string condition $\ell_1-\ell_2 = -\tfrac{9i}{2}N\hbar$.
\end{itemize}

\paragraph{One-loop ($O(\hbar^0)$), Toda Fermi-edge, and defect-fusion synthesis (Routes A, D, E; Phases 2b--2h).}
\begin{itemize}
\item \emph{Candidate Gamma S-matrix tested and ruled out at $O(\hbar^0)$:} A natural Gamma-function ansatz $S \propto \frac{\Gamma(1\pm i\eta)}{\Gamma(1\mp i\eta)}$ in $\eta = \frac{2}{9}(\ell_1-\ell_2)/\hbar$ (\texttt{smatrix\_ansatz.pdf}, \Ld) matches the semiclassical limit, unitarity, crossing, and the Bogdan--Charkina breather poles, but is \emph{not exact}: it is ruled out at $O(\hbar^0)$ by the quantum Toda Fermi-edge limit (\texttt{toda\_edge.pdf}), by the direct one-loop DHN calculation on the good equation (\texttt{oneloop\_phase.pdf}), and by KP-triangle defect fusion (\texttt{oneloop\_defect\_phase.pdf}).
\item \emph{Exact one-loop Legendre $\chi_2$ phase and canonical soliton momentum:} On the good equation, phonon scattering on $N$ solitons and Corrigan--Zambon defects is reflectionless (\texttt{oneloop\_phase.pdf}, \texttt{oneloop\_defect\_phase.pdf}). While $\delta$ at $O(\hbar^{-1})$ is elementary, the one-loop ($O(\hbar^0)$) phase evaluates in closed form to a 4-vertex KP coboundary of Legendre's dilogarithm $\chi_2(z) = \tfrac12(\operatorname{Li}_2(z)-\operatorname{Li}_2(-z))$ plus a logarithmic coboundary. In the co-moving one-loop effective Lagrangian $L_1(v) = v\Delta P_{\rm mode}-\Delta E_1^{\rm ren}$, the light root $\phi_3=\alpha$ drops out of $dL_1/d\alpha$ identically, yielding an $\alpha\ln\alpha$-free canonical collective-coordinate momentum $\Delta P_{\rm can}(\alpha,\epsilon) = -\frac{4\epsilon}{\pi}\alpha + O(\alpha^3)$ satisfying $d\Delta E_{\rm can} = v\,d\Delta P_{\rm can}$ (\texttt{charge\_frame\_master.pdf}).
\item \emph{Exact analytical match to the quantum Toda Fermi-edge limit:} In the quantum Toda Fermi-edge limit (\texttt{toda\_edge\_ba.py}, \texttt{matching.pdf}, \texttt{charge\_frame\_master.pdf}), measuring dressed momentum relative to the quantum Fermi edge absorbs the $\ln\hbar_T$ divergence into $q$, the classical frame term $+4\sqrt{u_1 u_2}=-M_1\Pi_2/\hbar_T$ cancels the $-2/r$ Darboux singularity in $\partial\Phi_{\rm cl}/\partial u_2 = -2\ln\frac{1-r}{1+r}$, and the overtaking amplitude-ordering reversal ($v_{\rm good}=\cos\alpha$ vs $v_{\rm bad}=\cosh\beta$) flips $H_{\rm bad}(r)=-H_{\rm good}(r)=+\frac{2}{\pi}\ln r\ln\frac{1-r}{1+r}$, leaving the exact $Y\to\infty$ residual $\lim_{Y\to\infty} R_0(Y,r) = -\frac{4}{\pi}(\chi_2(r)-r)$, which matches the numerical Phase~2f table to $<7.1\times 10^{-5}$.
\item \emph{Universal KP-Coboundary Master Function:} For any odd kernel $K(z)=-K(-z)$, the soliton--defect phase $\mathcal{I}_{SD}[K]$ and two-soliton phase $\mathcal{I}_{SS}[K]$ are the first and second KP-chord differences of a $2\times 2$ dipole--dipole master function $\mathcal{F}_K(p,q)$ on the KP root space (\texttt{charge\_frame\_master.pdf}). We end with an updated ranked list of next steps towards the higher-loop / all-loop kernel $K(z;\hbar)$.
\end{itemize}
\end{abstract}
\tableofcontents

\section{Setting}
Throughout, $v=2(\log\tau)_{xx}$ and $k=\sinh\beta$. A soliton has velocity $v=\epsilon\cosh\beta$, with $\epsilon=+1$ for right-movers and $\epsilon=-1$ for left-movers. The potential is $v=\phi_x$ and the Lagrangian is
\[
\mathcal L=\tfrac12\phi_t^2-\tfrac12\phi_x^2+\tfrac12\phi_{xx}^2+\phi_x^3 ,
\]
with $\pi=\phi_t$ and canonical symplectic form $\Omega=\int\delta\pi\wedge\delta\phi$.

On one soliton the charges are:
\begin{itemize}
\item mass $M=\int v=2k$;
\item $\phi$-shift charge $\Pi=\int\pi=-vM$, conserved because $\phi\to\phi+c$ is a symmetry;
\item momentum $P=-\int\pi\phi_x=\tfrac23\epsilon k^3\cosh\beta$;
\item energy $E=\tfrac2{15}k^3(5+4k^2)$.
\end{itemize}
These satisfy $dE/dP=v$. \La~\cite{BC}/\Lc

The two-soliton constants are
\[
A_{\rm ov}=\frac{1+3\coth^2\frac{\beta_1+\beta_2}{2}}{1+3\coth^2\frac{\beta_1-\beta_2}{2}},\qquad A_{\rm ho}=\frac{1+3\tanh^2\frac{\beta_1-\beta_2}{2}}{1+3\tanh^2\frac{\beta_1+\beta_2}{2}} ,
\]
both in $(0,1)$, and the position shifts are $\Delta x_1=-k_1^{-1}\log A$, $\Delta x_2=+k_2^{-1}\log A$ (soliton 1 starts on the left). \Lc~(Phase~1)

\paragraph{What a soliton S-matrix would mean here.} The model has no boost symmetry and is gapless (phonons with $\omega\simeq|k|$). The natural object is therefore the non-relativistic two-body S-matrix $S(P_1,P_2)$ of soliton asymptotic states, in the sense of Faddeev--Korepin \cite{FK}. Classical scattering is purely transmissive and diagonal, so $S$ is a phase; for two species (left and right movers) $S$ is diagonal.

For diagonal scattering the Yang--Baxter equation is automatic. The non-trivial requirements are:
\begin{itemize}
\item unitarity;
\item the correct semiclassical limit;
\item analyticity and a crossing substitute;
\item a bound-state (breather) spectrum consistent with poles.
\end{itemize}

%=====================================================================
\section{Route A: semiclassical quantisation}\label{sec:A}

\subsection{The symplectic structure of soliton scattering}

\begin{proposition}[\Lc]\label{prop:Omega}
Consider two well-separated solitons, with $\phi$ normalised by $\phi(-\infty)=c$. Up to exponentially small terms,
\[
\Omega=d\Pi_{\rm tot}\wedge dc+\sum_{j}dP_j\wedge dX_j+\big[X_i>X_j\big]\,d\Pi_i\wedge dM_j .
\]
\end{proposition}

The argument is as follows. $\pi_i$ is localised at $X_i$, and there $\phi_j$ equals $0$ or $M_j$ according to whether soliton $i$ is to the left or right of soliton $j$. The self-term $dP\wedge dX$ was derived in Phase~1. In sine-Gordon the analogous cross term vanishes, because the kink height $2\pi/\beta$ is fixed. Here $M=2k$ varies with the soliton parameter.

\begin{proposition}[\Lc, \script{symplectic\_scattering.py}]\label{prop:symp}
The classical scattering map $X_j\mapsto X_j+\Delta x_j$ (with $k_j$ and $c$ fixed) preserves $\Omega$, i.e.
\[
\sum_jdP_j\wedge d\Delta x_j=(\text{cross term})_{\rm in}-(\text{cross term})_{\rm out}.
\]
This was verified to 40 digits for overtaking and head-on collisions at five parameter points each. It was also verified for the \emph{good} equation in its regular region (\script{consistency\_checks.py}). Without the cross terms the identity fails by $O(1)$ amounts; for example, the left side equals $19.79$ at $\beta=(0.9,0.3)$.
\end{proposition}

This is a non-trivial confirmation of the classical picture. It also shows that a semiclassical S-matrix built from the naive rule ``$\Delta x_j=-\partial\delta/\partial P_j$'' does not exist: the shifts are not curl-free in the canonical momenta.

\subsection{The semiclassical phase}

In each ordering region we choose Darboux coordinates by the rule
\[
X_j'=X_j+\gamma_j\sum_{i\text{ left of }j}M_i,\qquad \gamma_j=\frac{d\Pi_j/dk_j}{dP_j/dk_j}.
\]
The mirrored rule (normalising $\phi$ at $+\infty$) gives the same in--out differences. The shifts $\Delta x_j'$ are then curl-free. Their potential is the semiclassical phase $\delta$, with $\Delta x_j'=-\partial\delta/\partial P_j$ and $S\simeq\exp(i\delta/\hbar)$, by stationary phase on two-particle wave packets \cite{FK}.

\begin{proposition}[\Lc, \script{semiclassical\_phase.py}, \script{consistency\_checks.py}]\label{prop:delta}
\begin{enumerate}
\item $\dfrac{dP}{d\beta}=\tfrac23\,k\,\sinh3\beta$ for the bad equation. For the good equation, $\dfrac{dP}{d\alpha}=\tfrac23\,k\sin3\alpha$ with $k=\sin\alpha$; this vanishes at the cusp $\alpha=\pi/3$ found in Phase~1.
\item For both overtaking and head-on collisions, and both equations,
\begin{equation}\label{eq:delta}
\boxed{\ \delta(1,2)=\tfrac29\,(\ell_1-\ell_2)\,\log A_{12}+\tfrac23\,M_1M_2\,(v_1-v_2)\ }
\end{equation}
with $\ell=\epsilon\cosh3\beta$ (bad) or $\ell=-\epsilon\cos3\alpha$ (good). The remainder was first found by symbolic integration, then identified as $\tfrac23M_1M_2(v_1-v_2)=-\tfrac23(\Pi_1M_2-\Pi_2M_1)$ by exact matching.
\item \emph{Unitarity:} $\delta(1,2)=-\delta(2,1)$.
\item \emph{Crossing-type relation:} $\delta_{\rm ov}(\beta_1,\beta_2+i\pi)=-\delta_{\rm ho}(\beta_1,\beta_2)$. The shift $\beta\to\beta+i\pi$ maps a right-mover to a left-mover with $M\to-M$.
\end{enumerate}
\end{proposition}

\begin{remark}
\begin{itemize}
\item $\ell$ is proportional to the spectral parameter of the third-order Lax operator. In the uniformisation of \cite{CL-WW}, $\lambda=\tfrac12(k^3+k^{-3})$, with $k^3$ playing the role of $e^{3\beta}$.
\item The dynamical part of $\delta$ is (difference of spectral parameters)$\times\log A$.
\item The second part is a ``mutual statistics'' phase: an Aharonov--Bohm-like phase $\propto\Pi_1M_2-\Pi_2M_1$ between objects carrying $\phi$-shift charge and $\phi$-winding. \Ld\ (interpretation)
\end{itemize}
\end{remark}

\begin{remark}[KdV limit, \Lc]
For small amplitudes the topological term is of the \emph{same order} as the KdV-type term $2\beta_1\log A$ in $\partial\delta/\partial\beta_1$. Numerically, their ratio tends to $0.564$ for $\beta_1/\beta_2=1.7$. The heuristic of Bonnemain and Doyon \cite{BD} identifies the quantum phase derivative with $P'_{\rm GHD}\Delta$, with $P_{\rm GHD}=\eta^2/2$. That identification reproduces our dynamical term but \emph{misses the topological term}. The classical generalised hydrodynamics, which uses only $\Delta x$, is unaffected. \Lc
\end{remark}

\subsection{Bohr--Sommerfeld: breathers as spectral-parameter strings}

Bogdan and Charkina quantise the Tajiri--Murakami breather of the bad equation by Bohr--Sommerfeld, $I_{\rm br}=N\hbar$, with an independently computed adiabatic invariant $I_{\rm br}$ \La~\cite[(64)--(65)]{BC}.

\begin{proposition}[\Lc]\label{prop:breather}
Set $\beta_{1,2}=a\pm ib$, which is Bogdan--Charkina's parametrisation $\kappa\pm ik=\sinh(a\pm ib)$. Then:
\begin{enumerate}
\item $A_{12}$ equals their $a_{12}$ \cite[(33)]{BC}. Regularity $A>1$ holds if and only if $b>\pi/3$, which is their existence boundary.
\item $i\cdot\tfrac29(\ell_1-\ell_2)=I_{\rm br}$ \emph{exactly}, with $I_{\rm br}=-\tfrac49\sinh3a\sin3b$, in agreement with \cite[(64)]{BC}.
\end{enumerate}
Hence the coefficient of $\log A$ in~\eqref{eq:delta} is the breather action. Bohr--Sommerfeld quantisation reads
\begin{equation}\label{eq:string}
\ell_1-\ell_2=-\tfrac{9i}{2}\,N\hbar,\qquad N\in\mathbb N ,
\end{equation}
which is a string condition in the spectral parameter $\ell$.
\end{proposition}

This is the strongest cross-check in this report. It ties our two-soliton phase, built from Hirota shifts and the canonical symplectic form, to an independent quasiclassical calculation in the literature. Item~2 was checked at three parameter points to $10^{-16}$.

Equation~\eqref{eq:string} is the form in which bound states appear in Bethe-ansatz solvable models. It suggests that $\ell$ is the variable in which an exact S-matrix has poles: $\ell_1-\ell_2\in-\tfrac{9i}2\hbar\mathbb N$, possibly with $\hbar$ renormalised. \Ld

\subsection{What Route A gives and subsequent updates (Phases 2c, 2e, 2h)}

\begin{itemize}
\item \emph{Obtained at tree level ($O(\hbar^{-1})$).} The leading semiclassical phase~\eqref{eq:delta}, which is unitary, crossing-type covariant and consistent with breather quantisation. Also an asymptotic Bethe-ansatz quantisation of $N$ solitons on a ring, $P_j\Lambda+\sum_{k\ne j}\delta(j,k)=2\pi\hbar n_j$. \Ld: the twisted boundary conditions forced by $\phi$-winding on a ring have not been worked out. Unlike in the relativistic Toda case \cite{Hol93,Johnson}, the tree-level phase $\delta$ is \emph{elementary}: no dilogarithm appears at $O(\hbar^{-1})$, because $dP\propto k\sinh3\beta\,d\beta$ and integration by parts closes.
\item \emph{Update (Phase 2c, \texttt{smatrix\_ansatz.pdf}): Candidate Gamma ansatz.} Setting $\eta=\tfrac29(\ell_1-\ell_2)/\hbar$ and $R(x)=\Gamma(1+ix)/\Gamma(1-ix)$, a candidate exact S-matrix was proposed \Ld: $S=R(\eta)(e^2 A/\eta^2)^{i\eta}e^{iT/\hbar}$ (overtaking) and $S=R(\eta)^{-1}(A\eta^2/e^2)^{i\eta}e^{iT/\hbar}$ (head-on). Its semiclassical expansion gives $\arg S - \delta/\hbar \to \pm\tfrac\pi2\operatorname{sgn}\eta + O(\hbar)$ and its poles match the Bogdan--Charkina breather levels~\eqref{eq:string}.
\item \emph{Update (Phase 2d, \texttt{oneloop.pdf}): Bad-equation one-loop DHN analysis.} For the bad equation, modes with $|q|>1$ are exponentially unstable ($|A(\text{sol},\text{ph})|\ne 1$), so the DHN sum is well posed only on the Charlier--Lenells invariant stable subspace $|q|<1$ (Scheme~A). There the breather one-loop energy and momentum are holomorphic continuations of the constituent soliton values, so the string condition~\eqref{eq:string} receives no $O(\hbar^0)$ shift at fixed rapidity; however, a reversed-group-velocity mode at $q^*=-\sqrt3/2$ is captured by slow solitons, making Scheme~A differ from the Toda continuum limit at $\beta\to 0$. \Lc
\item \emph{Update (Phase 2e, \texttt{oneloop\_phase.pdf}): Good-equation one-loop two-soliton phase.} For the well-posed good equation ($k=\sin\alpha$, $v=\epsilon\cos\alpha$), the $4\times 4$ linearised operator is reflectionless for all $q\in\mathbb R$, with a 3-root continuum phase shift $\theta_S'(q;\alpha,\epsilon)=\sum_{r=1}^3 c_r\Theta'(q;\phi_r^{(\epsilon)})$, weights $(c_1,c_2,c_3)=(1,1,-2)$, and angles $(\phi_1,\phi_2,\phi_3)=(\alpha+2\pi/3,\alpha-2\pi/3,\alpha)$. By Levinson's theorem, the mode count jumps at the cusp $\alpha=\pi/3$ where a heavy-branch instability appears. On the two-soliton background, the Faddeev--Korepin one-loop phase evaluates in closed form to a $3\times 3$ sum over $L(z)=-2\chi_2(z)-\ln|z|\ln\big|\frac{1-z}{1+z}\big|$, where $\chi_2(z)=\tfrac12(\operatorname{Li}_2(z)-\operatorname{Li}_2(-z))$ is Legendre's dilogarithm! Thus \textbf{dilogarithms enter at one loop ($O(\hbar^0)$) rather than at tree level ($O(\hbar^{-1})$)}, ruling out both the continued Gamma candidate~(A) and the pure exponential~(B) as exact. \Lc
\item \emph{Update (Phase 2h, \texttt{charge\_frame\_master.pdf}): One-loop soliton charge audit.} Because the Boussinesq soliton profile $u_0(x-vt;\alpha)$ changes width $k=\sin\alpha=\sqrt{1-v^2}$ with velocity $v=\epsilon\cos\alpha$, the fixed-profile mode-sum charges $(\Delta E_1^{\rm ren},\Delta P_{\rm mode})$ have a non-zero velocity anomaly $\mathcal A_v(\alpha)=d\Delta E_1^{\rm ren}/d\alpha - v\,d\Delta P_{\rm mode}/d\alpha\ne 0$. In the co-moving one-loop effective Lagrangian $L_1(v)=v\Delta P_{\rm mode}-\Delta E_1^{\rm ren}$, the total weight $\cos(2\phi_r)-\cos(\alpha+\phi_r)$ vanishes identically on the light root $\phi_3=\alpha$, cancelling all $\ln\alpha$ singularities to all orders and giving the canonical collective-coordinate momentum \Lc
\[
\Delta P_{\rm can}(\alpha,\epsilon)=\frac{\partial L_1}{\partial v}=-\epsilon\frac{\sqrt3\cos\alpha}{\pi}\ln\left|\frac{\tan\frac{\alpha+2\pi/3}{2}}{\tan\frac{\alpha-2\pi/3}{2}}\right|=-\frac{4\epsilon}{\pi}\alpha+\frac{8\epsilon}{9\pi}\alpha^3-\frac{8\epsilon}{45\pi}\alpha^5+O(\alpha^7),
\]
which satisfies $d\Delta E_{\rm can}/dv = v\,d\Delta P_{\rm can}/dv$ identically.
\end{itemize}

%=====================================================================
\section{Route B: quantum inverse scattering on the line}

\emph{What would have to be true.} The Lax matrix $U(x,\lambda)$ would need an ultralocal $r$-matrix Poisson algebra, so that it can be regularised on a lattice and quantised à la Faddeev--Korepin--Sklyanin. Then a Bethe ansatz over a pseudo-vacuum with soliton excitations would give $S$.

\begin{proposition}[\Lc, \script{routes\_BC\_checks.py}]
In canonical variables the Lax data of Phase~1 are $q=\tfrac32\phi_x-\tfrac14$ and $r=\tfrac12q_x+\tfrac{\sqrt3}4\pi$ (good-equation real form; the bad one differs by $i$'s). Their brackets are
\[
\{q(x),r(y)\}=\tfrac{3\sqrt3}{8}\,\delta'(x-y),\qquad\{q,q\}=\{r,r\}=0 .
\]
So the companion-matrix Lax connection is \emph{non-ultralocal}.
\end{proposition}

\emph{Where it stalls.} There are two known ways out, and neither yet gives a theory on the line.
\begin{enumerate}
\item Maillet's $r/s$ formalism with lattice regularisation \cite{Maillet}. No construction exists for Boussinesq.
\item The Miura (free-field) map. This is the route of Bazhanov--Hibberd--Khoroshkin \cite{BHK} and Hikami \cite{Hikami}. It yields $U_q(\widehat{\mathfrak{sl}}_3)$ T- and Q-operators, but in the chiral conformal setting on a circle, with no particle sector.
\end{enumerate}

Dimakis et al.'s Yang--Baxter maps \cite{DMHC} concern polarisations in the matrix case and are vacuous for the scalar equation. The ``beyond Yang--Baxter'' remark there is therefore not an obstruction here. For diagonal scattering, Yang--Baxter holds trivially. \Lc

%=====================================================================
\section{Route C: $W_3$, sl$_3$ affine Toda and the $U_q(\widehat{\mathfrak{sl}}_3)$ R-matrix}

\emph{What would have to be true.} Boussinesq solitons would have to be identifiable as a limit of $a_2^{(1)}$ (complex) affine Toda solitons, whose S-matrices are the $U_q(\widehat{\mathfrak{sl}}_3)$ intertwiners of Hollowood \cite{Hol93} and Gandenberger \cite{Gand}.

\begin{proposition}[\Lc, \script{routes\_BC\_checks.py}]\label{prop:fact}
Take Hollowood's $a_2^{(1)}$ interaction constant \cite[(1.19)]{Hol92}, $A^{(12)}(\theta)=\frac{2\cosh\theta+1}{2(1+\cosh\theta)}$, between the two fundamental species. With $S=\beta_1+\beta_2$ and $D=\beta_1-\beta_2$,
\[
A^{\rm bad}_{\rm ov}(\beta_1,\beta_2)=\frac{A^{(12)}(S)}{A^{(12)}(D)}\cdot\frac{\tanh^2(D/2)}{\tanh^2(S/2)},\qquad A^{\rm bad}_{\rm ho}(\beta_1,\beta_2)=\frac1{A^{\rm bad}_{\rm ov}(\beta_1,\beta_2+i\pi)} .
\]
Here $\tanh^2(\theta/2)$ is the sine-Gordon soliton--soliton constant. Equivalently,
\[
A_{\rm ov}=\frac{G(S/2)}{G(D/2)},\qquad G(y)=\frac{16\prod_{a=1,2}\sinh\frac{y+ia\pi/3}{2}\sinh\frac{y-ia\pi/3}{2}}{\sinh^2y}.
\]
\end{proposition}

So the Boussinesq interaction is built from exactly the $\mathfrak{sl}_3$ blocks, but in the form ``$F(\beta_1-\beta_2)/F(\beta_1+\beta_2)$''. This is the structure of scattering off a particle \emph{and its mirror image}. It suggests that a Boussinesq S-matrix might take the form
\[
S(\beta_1,\beta_2)=\mathcal S(\beta_1-\beta_2)\,\mathcal S(\beta_1+\beta_2+i\pi)^{\pm1}
\]
with $\mathcal S$ built from $U_q(\widehat{\mathfrak{sl}}_3)$ and sine-Gordon factors. Such a form is reminiscent of boundary (reflection) amplitudes. That would connect to the ``boundary affine Toda'' structure in Bazhanov--Hibberd--Khoroshkin \cite{BHK}. \Ld

The obstruction is that the semiclassical phase~\eqref{eq:delta} is not of this product form. Its weight $\tfrac29(\ell_1-\ell_2)$ is not a sum of functions of $S$ and $D$, and its topological part has no Toda analogue.

Further obstructions:
\begin{itemize}
\item Toda solitons are relativistic with fixed masses $m_a$. In the Toda constant the mass term $4m^2\sin^2(\pi a/3)$ appears where Boussinesq has $K^4$, so the species label is ``replaced'' by the wavenumber. There is no limit of fixed-species Toda solitons that yields Boussinesq solitons. \Lc
\item Takács and Watts show that $a_2^{(1)}$ soliton--breather scattering is flawed both classically and quantum mechanically \Lb~\cite{TW}.
\end{itemize}
\emph{Route C stalls}: the structural similarity is precise, but no dictionary is available. What is lost in the non-relativistic reduction is the difference form, and with it the standard crossing.

%=====================================================================
\section{Route D: discrete models and the quantum Toda Fermi-edge limit}

\emph{(i) Quantum discrete Boussinesq (Hikami \cite{Hikami}).} This is the Belov--Chaltikian lattice $W_3$ chain, periodic, with a Baxter $Q$-operator. A soliton/particle sector would require a thermodynamic limit with an identified ground state and excitations. That has not been done in the literature, and we did not attempt it. \emph{Stalls immediately}; it is a research project in itself.

\emph{(ii) Quantum Toda chain.} This is our recommended route for the small-amplitude corner. The required ingredients are:
\begin{enumerate}
\item The exact quantum Toda chain \La~\cite{Gutzwiller,Sklyanin} and its asymptotic Bethe ansatz \La/\Lb~\cite{Sutherland,Mertens,SS}. Its particle excitations become solitons in the classical limit \Lb~\cite{Mertens,PM}.
\item The dressed particle--particle S-matrix in the finite-density ground state.
\item Its classical limit should reproduce the Toda soliton phase.
\item The small-amplitude scaling limit should give~\eqref{eq:delta}.
\end{enumerate}

\paragraph{Update (Phase 2b, \texttt{toda\_route.pdf}).} Steps 2--4 were carried out in the classical limit, and step 2 at finite $\hbar$ to $O(\hbar)$. \Lc
\begin{itemize}
\item The classical limit of the Bethe ansatz reproduces the Toda soliton energies, momenta and exact two-soliton shifts.
\item Its two-body phases carry exactly the predicted topological cross terms.
\item The label-frame phase is $\delta_T=(p_1-p_2)\log A_T-(\Pi_1M_2-\Pi_2M_1)$.
\item In the scaling limit, $\delta_T\to3b\,\delta_{\rm bad}$ for overtaking \emph{and} head-on collisions. This independently confirms~\eqref{eq:delta}, including its topological term, from an exactly solvable quantum model.
\end{itemize}

\paragraph{Update (Phases 2d, 2f, 2h, \texttt{toda\_edge.pdf}, \texttt{matching.pdf}, \texttt{charge\_frame\_master.pdf}): Fermi-edge double-scaling limit and exact analytical resolution.}
Setting $y_j = \beta_j^2/\hbar_B = b\kappa_j^2/\hbar_T$ fixed as $\hbar_T=3b\hbar_B\to 0$ (\script{toda\_edge\_ba.py}, \script{matching\_toda.py}): \Lc
\begin{itemize}
\item \emph{Head-on channel:} The exact Fermi-edge phase equals $\delta/\hbar_B$ with vanishing $O(\hbar^0)$ residual ($<10^{-6}$ as $y_j\to\infty$), matching the small-amplitude head-on continuum limit $I_{\rm head}(\alpha_1,\alpha_2)=O(\alpha_1\alpha_2\ln(\alpha_1\alpha_2))\to 0$.
\item \emph{Overtaking channel ($y_1=Y > y_2=r^2 Y$, $r\in(0,1)$):} The Phase~2c Gamma candidate matches the Fermi-edge phase to $\lesssim 10^{-3}$ for $\eta\gtrsim 2.5$ but leaves a genuine two-body discrepancy of up to $\sim 0.1$ near threshold (\texttt{toda\_edge.pdf}). In the semiclassical limit $Y\to\infty$ after $\hbar_T\to 0$, the Toda residual $R_0(Y,r)$ approaches a smooth non-zero $O(\hbar^0)$ function $c_0(r)$ (\texttt{matching.pdf}, \texttt{matching\_tables.tex}).
\item \emph{Exact analytical resolution of $R_0(Y,r)$ (\texttt{charge\_frame\_master.pdf}):} Measuring dressed momentum relative to the quantum Fermi edge $P_{\rm dr}(p_j)-P_{\rm dr}(q)=P_n(\kappa_j)/a$ absorbs the $\ln\hbar_T$ ($\ln\lambda$) infrared divergence into $q$, shifting $u_j:=\lim_{\hbar_T\to 0}(p_j-q)/\hbar_T = y_j - \frac{1}{2\pi}\ln y_j + c_{\rm edge} + O(y_j^{-1})$. In the classical edge Bethe phase $\Phi_{\rm cl}(u_1,u_2)=2(u_1-u_2)\ln\frac{1-r}{1+r}+4\sqrt{u_1 u_2}$, the topological frame term $+4\sqrt{u_1 u_2}=-M_1\Pi_2/\hbar_T$ cancels the $-2/r$ Darboux singularity in $\partial\Phi_{\rm cl}/\partial u_2 = -2\ln\frac{1-r}{1+r}$ identically. Furthermore, because $v_{\rm good}(\alpha)=\cos\alpha$ decreases with amplitude while $v_{\rm bad}(\beta)=\cosh\beta$ increases with amplitude, overtaking collisions ($v_1>v_2$) reverse the amplitude ordering ($\alpha_1<\alpha_2$ vs $\beta_1>\beta_2$), flipping $H_{\rm bad}(r)=-H_{\rm good}(r)=+\frac{2}{\pi}\ln r\ln\frac{1-r}{1+r}$. Combining the quasi-momentum shift $d\Phi_{\rm cl}(Y,r)$ and the intrinsic loop shift $d\Phi_{\rm loop}(Y,r)$, both the $\pm 4r\delta u_1 \sim \mp\frac{2r}{\pi}\ln Y$ topological terms and the $\pm\frac{2}{\pi}\ln r\ln\frac{1-r}{1+r}$ logarithmic terms cancel identically, yielding the exact $Y\to\infty$ limit \Lc
\begin{equation}\label{eq:R0_exact}
\boxed{\ \lim_{Y\to\infty} R_0(Y,r) = -\frac{4}{\pi}\big(\chi_2(r) - r\big) = -\frac{4}{\pi}\sum_{n=1}^\infty \frac{r^{2n+1}}{(2n+1)^2}\ }
\end{equation}
(where $+\frac{4r}{\pi}$ is the Fermi-edge reference subtraction at $r\to 0$), which matches the numerical Phase~2f intercepts ($r=0.4,0.5,0.6$) to $1.2\text{--}7.0\times 10^{-5}$ and the $r$-slope to $2.9\times 10^{-4}$ (\script{charge\_and\_frame\_audit.py}).
\end{itemize}

%=====================================================================
\section{Route E: integrable defects and the Universal KP-Coboundary Master Function}

Corrigan and Zambon (CZ) constructed integrable type-I Bäcklund defects for the Boussinesq equation \La~\cite{CZ}. Each defect carries a single KP root $q_D=-s/(2\sqrt3)$ ($s=\eta\sigma$, $\eta=\pm 1$), whereas a soliton carries a KP chord $(p_a,p_b)$ connecting two roots of the KP cubic $4q^3-q-\ell/(3\sqrt3)=0$.

\paragraph{Semiclassical soliton--defect transmission (Phase 2c, \texttt{defect\_semiclassics.pdf}).} \Lc
\begin{itemize}
\item The classical soliton--defect transmission phase is $\delta_D = \frac29(\ell-\ell_D)\ln|z_t| - \frac23(\Pi M_D - \Pi_D M)$ with $z_t=((p_b-q_D)/(p_a-q_D))^{\epsilon\eta}$ and $\ell_D=(3s-s^3)/2$. The static defect adds no ordering-dependent cross terms; two-soliton + defect scattering is symplectic and factorises.
\item A defect vacuum $q_D$ is linearly stable if and only if $s = \eta\sigma < -1$ ($q_D > 1/(2\sqrt3)$), which selects the unique ground state in each KP triple $\{q_1,q_2,q_3\}$.
\item Along any edge $q_i\to q_{i+1}$ of a KP triangle where a soliton $S_i$ on $\{q_i,q_{i+1}\}$ is absorbed by the defect, the tree-level phases satisfy exact defect fusion $\delta_D(q_{i+1})(3)-\delta_D(q_i)(3)=\rho_{{\rm tree},i}[\delta(3,S_i)+M_{S_i}\Pi_3]$.
\end{itemize}

\paragraph{One-loop soliton--defect phase and KP-triangle closure (Phase 2g, \texttt{oneloop\_defect\_phase.pdf}).} \Lc
\begin{itemize}
\item Phonon scattering on $N$-soliton + CZ defect backgrounds is reflectionless ($R(q)=0$) with factorised transmission amplitude $T_{\rm tot}(q)=\big(\prod_{j=1}^N A(j,{\rm ph}(q))\big)z_{\rm ph}(q;q_D,\eta)$.
\item Setting $q_D=-(\cos\gamma)/\sqrt3$, the defect phase-shift derivative is a \emph{2-root $(-\eta,+\eta)$ dipole} $\theta_D'(q;s,\eta)=\operatorname{Re}\sum_{s'=1}^2 c_{D,s'}\Theta'(q;\psi_{s'})$ at $\psi_{1,2}=\pm 2\pi/3-\gamma$, whose difference across $(p_a,p_b)$ reproduces the 3-root $(1,1,-2)$ soliton kernel $\theta_S'(q;\alpha,\epsilon)$.
\item The $4\times 4$ defect-vacuum localised-mode determinant factors into two quadratics with closed-form eigenvalues $\omega_{1,2}(\sigma,\eta)=\frac{i}{8}[\eta\sigma\sqrt{4-\sigma^2}\pm\sqrt3(2-\sigma^2)]$, explaining the Levinson jumps at $s=-1$ and $s=2$.
\item Because the static defect has no translational zero mode ($\theta_D^{\rm phys}(0^\pm)=0$), $I_{SD}=J_{SD}^{\rm phys}$ evaluates in closed form to a $3\times 2$ Legendre $\chi_2$ sum. One-loop defect vacuum shifts $(\Delta E_D^{\rm ren},\Delta P_D)$ balance $(\Delta E_1^{\rm ren},\Delta P_{\rm mode})$ in soliton absorption, and one-loop defect fusion $I_{SD}(3;q_{i+1},\eta)-I_{SD}(3;q_i,\eta)=\rho_i I_{SS}(3,S_i)$ holds on every KP edge and closes around every KP triangle ($\sum_{i=1}^3 \rho_i I_{SS}(3,S_i)=0$), whereas the $+\tfrac\pi2\delta_{\epsilon_3,\epsilon_{S_i}}$ constant of $J_{SS}$ fails to close by $\pm\pi/2$.
\end{itemize}

\paragraph{Universal KP-Coboundary Master Function (Phase 2h, \texttt{charge\_frame\_master.pdf}).} \Lc
For \emph{any} odd ratio kernel $K(z)=-K(-z)$, the $2\times 2$ dipole--dipole \textbf{two-vertex master function} on the KP root space (\script{kp\_master\_and\_toda\_edge.py}),
\begin{equation}\label{eq:F_master}
\mathcal F_K(p,q) := \sum_{s_1=1}^2\sum_{s_2=1}^2 (-1)^{s_1+s_2}\operatorname{Re} K\!\left(\tan\frac{\psi_{s_1}(p)}{2}\cot\frac{\psi_{s_2}(q)}{2}\right),
\end{equation}
generates the soliton--defect functional as its first difference $\mathcal I_{SD}[K](\alpha,\epsilon;q_D,\eta)=\eta(\mathcal F_K(p_b,q_D)-\mathcal F_K(p_a,q_D))$ and the two-soliton functional as its 4-vertex second difference
\begin{equation}\label{eq:ISS_coboundary}
\boxed{\ \mathcal I_{SS}[K](1,2) = \mathcal F_K(p_{1,b},p_{2,b}) - \mathcal F_K(p_{1,a},p_{2,b}) - \mathcal F_K(p_{1,b},p_{2,a}) + \mathcal F_K(p_{1,a},p_{2,a})\ }
\end{equation}
guaranteeing defect fusion and KP-triangle closure for any odd $K(z)$. Splitting the one-loop kernel into $K_L=K_{\chi_2}+K_{\ln}$ with $K_{\chi_2}(z)=-\frac1\pi\chi_2(z)$ and $K_{\ln}(z)=-\frac{1}{2\pi}\ln|z|\ln\big|\frac{1-z}{1+z}\big|\operatorname{sgn}(\operatorname{Re}z)$ shows that at finite amplitude $\mathcal I_{SS}[K_{\ln}]$ is the logarithmic quasi-momentum coboundary (which cancels $H_{\rm bad}(r)$ in the small-amplitude overtaking corner), leaving the pure Legendre $\chi_2$ coboundary $\mathcal I_{SS}[K_{\chi_2}]\to -\frac{4}{\pi}\chi_2(r)$.

%=====================================================================
\section{Consistency checks across candidates and orders}

\begin{center}\small
\begin{tabular}{@{}p{4.1cm}p{7.8cm}p{2.6cm}@{}}\toprule
Condition & Result & Status\\\midrule
Unitarity & $\delta(1,2)=-\delta(2,1)$ and $I_{SS}(1,2)=-I_{SS}(2,1)$, exact & \Lc\ passed\\
Crossing (substitute) & $\delta_{\rm ov}(\beta_1,\beta_2+i\pi)=-\delta_{\rm ho}(\beta_1,\beta_2)$; $I_{SD}(\alpha,-1;s,\eta)=I_{SD}(\alpha,+1;-s,\eta)$ & \Lc\ passed\\
Yang--Baxter / KP-triangle fusion & Diagonal YBE trivial; tree-level $\delta$ and one-loop $I_{SS}$ close around all 12 KP defect triangles ($\sum_i\rho_i I_{SS}=0$); Gamma candidate's $\pm\pi/2$ constant fails closure by $\pm\pi/2$ & \Lc\ decisive\\
Semiclassical limit vs.\ time delays & $\Delta x_j'=-\partial\delta/\partial P_j$ by symplectic consistency (Prop.~\ref{prop:symp}); canonical one-loop momentum $\Delta P_{\rm can}$ obeys $d\Delta E_{\rm can}=v\,d\Delta P_{\rm can}$ with zero $\alpha\ln\alpha$ term & \Lc\ passed\\
Bound states / bootstrap & Breather existence boundary and Bogdan--Charkina action reproduced (Prop.~\ref{prop:breather}); one-loop string condition unshifted at fixed rapidity in Scheme~A (\texttt{oneloop.pdf}) & \Lc\ passed\\
Good equation (cusp \& resonances) & One-loop phase $I_{SS}$ finite and smooth across all 3 resonance lines ($A_{12}\to 0,\infty$); jumps by $\pm\pi$ at the cusp $\alpha=\pi/3$ where a heavy-branch instability appears (\texttt{oneloop\_phase.pdf}) & \Lc\ passed\\
Quantum Toda Fermi-edge limit & Tree-level $\delta_T\to 3b\,\delta_{\rm bad}$; $O(\hbar^0)$ head-on residual $=0$; $O(\hbar^0)$ overtaking residual $\lim_{Y\to\infty} R_0(Y,r)=-\frac{4}{\pi}(\chi_2(r)-r)$ matches Phase~2f table to $<7.1\times 10^{-5}$ (\texttt{charge\_frame\_master.pdf}) & \Lc\ passed\\
\bottomrule\end{tabular}\end{center}

%=====================================================================
\section{Conclusions and ranked open problems}

\paragraph{Conclusions (Phases 2a--2h).}
\begin{enumerate}
\item \emph{Tree-level ($O(\hbar^{-1})$) and one-loop ($O(\hbar^0)$) two-soliton scattering are completely solved and reconciled.} At $O(\hbar^{-1})$, the phase is~\eqref{eq:delta}, where the ordering-dependent topological cross term $d\Pi_i\wedge dM_j$ is essential. At $O(\hbar^0)$, the continuum phase is the 4-vertex KP coboundary~\eqref{eq:ISS_coboundary} of the dilogarithm kernel $K_L = K_{\chi_2}+K_{\ln}$, whose small-amplitude Fermi-edge limit reproduces the quantum Toda Bethe ansatz residual $-\frac{4}{\pi}(\chi_2(r)-r)$ to $O(10^{-5})$. \Lc
\item \emph{Both simple candidates (the Gamma ansatz and the pure exponential) are ruled out at $O(\hbar^0)$.} A single Gamma factor $\Gamma(1+i\eta)/\Gamma(1-i\eta)$ in $\eta=\frac29(\ell_1-\ell_2)/\hbar$ has only rational/Bernoulli asymptotic corrections in $\eta$ and cannot produce the dilogarithm $\chi_2(z)$ at $O(\hbar^0)$; moreover, its $\pm\pi/2$ constant violates KP-triangle defect fusion. \Lc
\item \emph{Every physical soliton observable is a KP-chord difference.} Both at $O(\hbar^{-1})$ ($\ln A_{12}$ as the 4-vertex cross-ratio of $\ln|p-q|$) and at $O(\hbar^0)$ ($\mathcal I_{SS}[K]$ as the 4-vertex second difference~\eqref{eq:ISS_coboundary} of the $2\times 2$ dipole--dipole master function $\mathcal F_K(p,q)$), two-soliton scattering is governed by a two-vertex master function on the KP root space. \Lc
\end{enumerate}

\paragraph{Updated ranked next steps (after Phase 2h).}
\begin{enumerate}
\item \textbf{Analytical $O(1/Y)$ (two-loop, $O(\hbar^1)$) expansion of the quantum Toda Fermi-edge Bethe ansatz.} Carry the analytical Fermi-edge expansion of \script{toda\_edge\_ba.py} one order further in $1/Y$ (combining the $O(\hbar/D)$ Stirling/Bernoulli tail of the Toda loggamma kernel, the $O(1/y_j)$ quasi-momentum shift, and the rank-1 Hessian of $\Phi_{\rm cl}$) to derive a closed-form formula for the two-loop coefficient $c_1(r)$ extracted numerically in \script{kp\_master\_and\_toda\_edge.py} ($c_1(0.4)=-0.00708$, $c_1(0.5)=-0.02099$, $c_1(0.6)=-0.04822$).
\item \textbf{All-loop KP master kernel $K(z;\hbar)$ and quantum dilogarithm test.} Use the exact tree-level ($\ln$), one-loop ($\chi_2$), and two-loop ($c_1(r)/Y$) data together with the Universal KP-Coboundary representation~\eqref{eq:ISS_coboundary} to test whether the all-loop master function $\mathcal F_K(p,q;\hbar)$ is given by a quantum dilogarithm (Faddeev--Volkov / Barnes double Gamma function) on the KP root space.
\item \textbf{Ring quantisation with $\phi$-winding and twisted boundary conditions.} Formulate the periodic $N$-soliton Bethe equations incorporating the topological cross term $-\frac23(\Pi_i M_j-\Pi_j M_i)$ and the one-loop phase $\mathcal I_{SS}[K_L]$.
\item \textbf{Non-ultralocal quantum inverse scattering} (Maillet) for the Boussinesq Lax pair (lowest priority).
\end{enumerate}

\appendix\section{Phase 2 Write-ups and Verification Scripts}
All write-ups are in \texttt{phase2/*.pdf} and all scripts are in \texttt{phase2/scripts/*.py} (plus \texttt{corrigan-zambon/defect\_*.py}); each script ends with \texttt{ALL CHECKS PASSED}.
\begin{itemize}
\item \textbf{Phase 2a (\texttt{report.pdf}):} \script{symplectic\_scattering.py}, \script{semiclassical\_phase.py}, \script{fit\_R.py}, \script{consistency\_checks.py}, \script{routes\_BC\_checks.py}.
\item \textbf{Phase 2b (\texttt{toda\_route.pdf}):} \script{toda\_bethe\_classical.py}, \script{toda\_bethe\_quantum.py}, \script{toda\_ba\_vs\_label.py}, \script{toda\_scaling\_fast.py}, \script{toda\_quantum\_correction.py}.
\item \textbf{Phase 2c (\texttt{smatrix\_ansatz.pdf}, \texttt{defect\_semiclassics.pdf}):} \script{smatrix\_ansatz.py}, \script{smatrix\_boot.py}, \script{defect\_smatrix\_good\_continuation.py}, \texttt{corrigan-zambon/defect\_*.py}.
\item \textbf{Phase 2d (\texttt{toda\_edge.pdf}, \texttt{oneloop.pdf}):} \script{toda\_edge\_ba.py}, \script{toda\_edge\_twoterm.py}, \script{oneloop\_step0\_wellposedness.py}, \script{oneloop\_phase\_shifts.py}, \script{oneloop\_soliton\_energy.py}, \script{oneloop\_breather.py}.
\item \textbf{Phase 2e (\texttt{oneloop\_phase.pdf}):} \script{oneloop\_good\_phonon\_scattering.py}, \script{oneloop\_good\_twosoliton\_phase.py}, \script{oneloop\_good\_limits\_cusp\_resonance.py}.
\item \textbf{Phase 2f (\texttt{matching.pdf}):} \script{matching\_conventions.py}, \script{matching\_toda.py}, \script{matching\_analysis.py}.
\item \textbf{Phase 2g (\texttt{oneloop\_defect\_phase.pdf}):} \script{oneloop\_defect\_scattering\_and\_stability.py}, \script{oneloop\_defect\_phase\_and\_fusion.py}.
\item \textbf{Phase 2h (\texttt{charge\_frame\_master.pdf}):} \script{charge\_and\_frame\_audit.py}, \script{kp\_master\_and\_toda\_edge.py}.
\end{itemize}

\emph{Failures along the way (Phase 2a; later phases record their own failures in each write-up).}
\begin{itemize}
\item The direct symbolic simplification of the symplectic identity timed out, so we switched to 40-digit numerics.
\item The good-equation phase first used the wrong sign of $\ell$, and the test failed until it was corrected.
\item A tolerance set below double precision caused spurious failures.
\item The KdV-limit comparison with Bonnemain--Doyon \emph{does not} agree. This is recorded as a finding, not a passed test.
\end{itemize}

\begin{thebibliography}{99}\small
\bibitem{BC} M.~M.~Bogdan, O.~V.~Charkina, arXiv:2502.14695, \emph{Low Temp. Phys.} 51 (2025).
\bibitem{BD} T.~Bonnemain, B.~Doyon, \emph{SciPost Phys.} 18 (2025) 075, arXiv:2402.08669.
\bibitem{BHK} V.~V.~Bazhanov, A.~N.~Hibberd, S.~M.~Khoroshkin, \emph{Nucl. Phys. B} 622 (2002) 475, arXiv:hep-th/0105177.
\bibitem{CL-WW} C.~Charlier, J.~Lenells, arXiv:2204.02365.
\bibitem{CZ} E.~Corrigan, C.~Zambon, \emph{J. Phys. A} 56 (2023) 405202, arXiv:2306.03772.
\bibitem{DMHC} A.~Dimakis, F.~Müller-Hoissen, X.-M.~Chen, \emph{Phys. Scr.} 94 (2019) 035206, arXiv:1805.09711.
\bibitem{FK} L.~D.~Faddeev, V.~E.~Korepin, Quantum theory of solitons, \emph{Phys. Rep.} 42 (1978) 1, doi:10.1016/0370-1573(78)90058-3.
\bibitem{Gand} G.~M.~Gandenberger, arXiv:hep-th/9703158.
\bibitem{Gutzwiller} M.~C.~Gutzwiller, \emph{Ann. Phys.} 133 (1981) 304.
\bibitem{Hikami} K.~Hikami, \emph{Nucl. Phys. B} 604 (2001) 580, arXiv:nlin/0102021.
\bibitem{Hol92} T.~J.~Hollowood, \emph{Nucl. Phys. B} 384 (1992) 523, arXiv:hep-th/9110010.
\bibitem{Hol93} T.~J.~Hollowood, \emph{IJMPA} 8 (1993) 947, arXiv:hep-th/9203076.
\bibitem{Johnson} P.~R.~Johnson, \emph{Nucl. Phys. B} 496 (1997) 505, arXiv:hep-th/9611117.
\bibitem{Maillet} J.-M.~Maillet, Kac--Moody algebra and extended Yang--Baxter relations in the O(N) non-linear sigma model, \emph{Phys. Lett. B} 162 (1985) 137, doi:10.1016/0370-2693(85)91075-5. [not read]
\bibitem{Mertens} F.~G.~Mertens, \emph{Z. Phys. B} 55 (1984) 353.
\bibitem{PM} M.~Pustilnik, K.~A.~Matveev, \emph{Phys. Rev. B} 92 (2015) 195146, arXiv:1507.05639.
\bibitem{Sklyanin} E.~K.~Sklyanin, \emph{Lect. Notes Phys.} 226 (1985) 196. [not machine-verified]
\bibitem{SS} R.~Siddharthan, B.~S.~Shastry, \emph{Phys. Rev. B} 55 (1997) 12196, arXiv:cond-mat/9611221.
\bibitem{Sutherland} B.~Sutherland, \emph{Rocky Mountain J. Math.} 8 (1978) 413.
\bibitem{TW} G.~Takács, G.~M.~T.~Watts, \emph{Nucl. Phys. B} 547 (1999) 538, arXiv:hep-th/9810006.
\end{thebibliography}
\end{document}
